% !PRINTED 265
\weberpage{265}
%<<CONT
taken with negative sign, are congruent to $r$ modulo $\frac{1}{2}(p-1)$.
One therefore obtains all values of this expression that are distinct
from one another and from $\pm 2$ by setting
\weq{10}{r = 1, 2, \cdots \frac{p-3}{2}}
Since here for every exponent $k$
\[
S^k = \begin{pmatrix} \sigma^k, & 0 \\ 0, & \sigma^{-k} \end{pmatrix}
\]
the degree of $S$, and hence also of every substitution of the
II\textsuperscript{nd} kind, is either $\frac{1}{2}(p-1)$ or a divisor of it, according
as $r$ is relatively prime to $\frac{1}{2}(p-1)$ or not.

To determine the number of these substitutions, we proceed
as above. If first $c = 0$, it follows that
\[
a = \gamma^{\pm r(p+1)}, \quad d = \gamma^{\mp r(p+1)},
\]
and $b$ can have $p$ values. The number $r$ takes the values (10),
and accordingly this yields $p\,(p-3)$ substitutions. If next
$c$ is different from zero, then $p\,(p-1)$ is the number of the
distinct assumptions for $a, c$. If $a$ is assumed, we can
determine $d$ from (9) and set
\[
b = -c^{-1} + a\,d\,c^{-1},
\]
from which we obtain $\frac{1}{2}\,p\,(p-1)(p-3)$ determinations,
and so, together with those counted for $c = 0$, $\frac{1}{2}\,p\,(p-3)(p+1)$.
This number too is still to be halved, so that we obtain
\[
\frac{1}{4}\,p\,(p-3)(p+1)
\]
substitutions of the II\textsuperscript{nd} kind.

III.~$(a+d)^2 - 4$ a quadratic non-residue of $p$. In this
case $\sigma$ is not real, but conjugate to its reciprocal value.
Hence by \S.\,67, (4) $\sigma^p = \sigma^{-1}$, and consequently
\[
\sigma^{\frac{p+1}{2}} = \pm 1.
\]
From this it follows that the degree of such a substitution
is $\frac{1}{2}(p+1)$ or a divisor of this number.
In this third case we set
\weq{11}{\begin{aligned}
&\sigma = \gamma^{r(p-1)}, \quad a+d = \gamma^{r(p-1)} + \gamma^{-r(p-1)}\\
&r = 1, 2, \ldots, \tfrac{1}{2}(p-1).
\end{aligned}}
Here the case $c = 0$ or $b = 0$ cannot occur, because
from $a\,d = 1$ it would follow that $a = \sigma^{\pm 1}$, $d = \sigma^{\mp 1}$
%CONT>>
